\section{A finite witness on individual words}
\label{route:section}

Corollary~\ref{tri:positivity} proves positivity of the elementary
coefficients. We now construct the partition assigned to each individual
word. The construction matches ordered packets of vertices with finite
wall histories and anchor words; the anchor multiplicities specify the
partition. Its local matchings use coefficients of elementary sections
and individual wall operators, without taking the final elementary
coefficient list as input.

Fix a triangular embedding with $D\geq\max\{1,n\}$, put $N=n$, and
let $u$ be the sum of the indicated vertex exponents.  Write $H_-$ for its all-negative chamber,
$M=|E(G)|$, and
\[
 R=u-Nd_1\in C^+,
 \qquad B=\{y:Nd_1\leq y\leq u\}.
\]
The order in this display is the free-cone order.  If
$R=\sum_{c\in C}R_c c$, set $h(R)=\sum_c R_c$.
The set $B$ has $\prod_c(R_c+1)$ elements.  Only wall roots $0<r\leq R$
can contribute to the coefficients considered here.  All chamber and
path choices in this section refer to this finite truncation.

\subsection{Decorated coefficients and their finite orders}
\label{route:decorations}

For $1\leq k\leq N$, let
\[
 E_k=e_k(Y)=\Theta_{d_1+\cdots+d_k},
 \qquad E_k^H=\sum_{a,e}c_k^H(a,e)v^eX_a.
\]
The identity is Theorem~\ref{tri:monomial}.  The coefficients
$c_k^H(a,e)$ are nonnegative integers by Theorem~\ref{wall:positivity}.
A decorated $k$-packet in $H$ is a triple $(a,e,t)$ with
$1\leq t\leq c_k^H(a,e)$.  The integer $e$ is its energy.  For an ordered
composition ${\boldsymbol k}=(k_1,\ldots,k_b)$ of $N$, a product state is
an ordered tuple of decorated packets.  Its exponent and energy are
\begin{equation}
 \label{route:product-energy}
 a=\sum_{i=1}^b a_i,
 \qquad
 E=\sum_{i=1}^b e_i+\sum_{i<j}\Omega(a_i,a_j).
\end{equation}
Thus these states count the coefficients of $P_{\boldsymbol k}
=E_{k_1}\cdots E_{k_b}$, with their multiplicities.  In $H_-$, a packet
is just an independent set of size $k$, with energy zero and index one.

If a wall operator $W$ is positive on $X_a$, a decorated jump from $a$
to $a+jp$ is a triple $(j,e,t)$ with
\begin{equation}
 \label{route:jump-decoration}
 j\geq0,\qquad
 1\leq t\leq [v^eX_{a+jp}]W(X_a).
\end{equation}
Here $p$ is the primitive positive lattice vector of the wall ray; an
unoccupied multiple simply has coefficient zero.  Energies in
Equation~\eqref{route:jump-decoration} are in centered torus notation and already
include the shift arising from multiplication by $X_a$.  The unique
zero jump is $(0,0,1)$.

\begin{lemma}[Effective coefficient slices]
\label{route:finite}
At fixed total exponent in $B$ and fixed total energy, every packet
decomposition and local jump decomposition used below is a finite set
that can be enumerated effectively.  The same holds for a sequence of at most
$h(R)$ nonzero total jumps.
\end{lemma}

\begin{proof}
The support of $E_k$ in every chamber is contained in $kd_1+C^+$.
Consequently every packet offset and every jump increment in a
decomposition with total exponent at most $u$ is at most $R$.
There are finitely many exponent tuples, wall directions, and truncated
chambers.  This assertion applies to underlying packet states as well
as to their states after jumps.

At each fixed exponent the coefficient of a section or wall jump lies
in $\Z((v^{-1}))$: it has a finite upper energy bound and finite
coefficients at individual energies.  For a fixed exponent tuple,
write its energy equation as
\[
 E=\kappa+e_1+\cdots+e_s,
\]
where $\kappa$ is the fixed sum of torus cross terms.  If the coefficient
of the $i$th piece has upper energy bound $U_i$, then
\begin{equation}
 \label{route:energy-bound}
 E-\kappa-\sum_{j\ne i}U_j\ \leq e_i\leq U_i.
\end{equation}
Thus each energy and each decoration ranges over a finite set.  For
bounded sequences there are again only finitely many exponent types,
so the same argument applies to all their pieces at once.

These are effective bounds.  The initial units have the rational
coefficients
\[
 [X_{jp}]D_p=
 \frac{v^{-j^2}}{\prod_{i=1}^j(1-v^{-2i})}.
\]
The triangular factorization and incoming equations of
Section~\ref{wall:section} determine every wall coefficient to the
required root order by finite rational arithmetic in $v$.  Transport
of the finite independent-set sums determines the needed $E_k^H$ in
the same way.  A rational function expanded at $v^{-1}=0$ has a
computable upper degree, and any finite interval of its coefficients
is computable by formal division.  Reducing numerator and denominator
first also decides whether it vanishes.  One may therefore compute
the bounds in Equation~\eqref{route:energy-bound}, enumerate all choices in
them, and compute their multiplicities.  No factorization into
positive units needs to be found in order to perform these
coefficient calculations; positivity justifies interpreting the
computed integers as multiplicities.
\end{proof}

We specify the finite choices used by these definitions. Order lattice
coordinates in the basis $d_1$, followed by the forward roots $C$ in
time order and then any added dual basis vectors.
For a nonnegative integer $a$,
let $b(a)$ be its binary expansion, with $b(0)=0$, and encode it as
$1^{|b(a)|}0b(a)$.  Encode a signed integer by a sign bit followed by
the code of its absolute value, a rational by its reduced numerator
and positive denominator, and a tuple by its length followed by its
component codes.  A fixed initial type tag distinguishes these
possibilities.  Sign vectors are tuples of signed integers.  Order
codes first by length and then lexicographically.  In every finite
set below, ``the same index'' means the same rank in this order.
Fix the triangular embedding by the same order: enumerate tuples giving
an integer $D\geq\max\{1,n\}$, rational phases in the order of
Equation~\eqref{tri:phase-order}, and $n$ distinct selected bridges,
and take the first tuple whose induced ordered graph is $G$.
Every candidate is checked by finitely many rational comparisons, and
Lemma~\ref{tri:embedding} proves that a candidate exists.  For empty
input use $D=1$ and no selected bridges.

All later generic rational choices mean the first code satisfying the
listed open inequalities and avoidance conditions.  These are decidable
finite systems of rational linear conditions.  In particular, for a
line $h+tV(y)$, two distinct wall hyperplanes with nonzero pairings
have simultaneous crossings precisely when
$h(r)\Omega(s,y)-h(s)\Omega(r,y)=0$.  A root $r$ with
$\Omega(r,y)=0$ instead requires $h(r)\ne0$.  Thus the path conditions
used below have the asserted finite form.  A nonempty open rational cell contains a rational
point, so the searches terminate.

\subsection{The local packet switch}
\label{route:local}

Let $W:I\longrightarrow II$ be a wall crossing on the positive ray
$\N p$.  Write its operator as $\operatorname{Ad}(g_p)^\varepsilon$,
where $\varepsilon=1$ for a crossing from positive to negative and
$\varepsilon=-1$ for the reverse crossing.  Assume that the target
exponent $y$ satisfies
\begin{equation}
 \label{route:positive-total}
 \varepsilon\Omega(p,y)\geq0.
\end{equation}
This is exactly the positive bend direction of
Proposition~\ref{wall:factors}, with the identity direction included
as in Lemma~\ref{wall:unit}.

The total pairing in Equation~\eqref{route:positive-total} can be
nonnegative even when an individual packet has negative pairing.
The switch must therefore use both directions of single-packet
transport. Its invariant is a tuple of packet states retained on
the side from which their jumps have nonnegative coefficients.

\begin{proposition}[Packet switch]
\label{route:switch}
There is a computable energy-preserving bijection
\begin{equation}
 \label{route:switch-equation}
 \left\{\text{product states of exponent $y$ in $II$}\right\}
 \longleftrightarrow
 \left\{\begin{array}{c}
  \text{product states of exponent $z$ in $I$, together with}\\
  \text{a decorated jump of $W$ from $z$ to $y$}
 \end{array}\right\}.
\end{equation}
On the right the energy is the state energy plus the jump energy.
If $\Omega(p,y)=0$, the total jump is always the unique zero jump.
\end{proposition}

\begin{proof}
For a packet with exponent $a$, call its type positive if
$\varepsilon\Omega(p,a)\geq0$, and negative otherwise.  Adding a
multiple of $p$ does not change this type.  In particular the types
are recoverable on both sides of Equation~\eqref{route:switch-equation}.

For a positive packet the section identity $E_k^{II}=W(E_k^I)$
expresses each decorated state in $II$ as a state in $I$ followed by
a positive jump of $W$.  More precisely, at fixed output exponent
and energy the two sets have the same cardinality, so identify them
by their finite ranks from Lemma~\ref{route:finite}.  For a negative
packet use instead $E_k^I=W^{-1}(E_k^{II})$: a state in $I$ is
identified with a state in $II$ followed by a positive jump of
$W^{-1}$.  Packets of zero pairing are fixed.  These identifications
have computable inverses.

Fix the retained state of each packet: in chamber $I$ for a positive
packet, and in chamber $II$ for a negative packet. Write $x_i$ for
its exponent, and keep its energy and decoration index fixed.
The two positive expansions to be matched are
\[
\begin{array}{c|c|c|c}
\text{packet type}&\text{retained chamber}&
 L_i\ \text{(side $II$)}&R_i\ \text{(side $I$)}\\ \hline
\text{positive}&I&W(X_{x_i})&X_{x_i}\\
\text{negative}&II&X_{x_i}&W^{-1}(X_{x_i})
\end{array}
\]
Via the single-packet bijections, a decorated term of $L_i$ together
with its retained state gives a packet in $II$; a decorated term of
$R_i$ together with that state gives a packet in $I$.
Each row satisfies $L_i=W(R_i)$, so the automorphism property gives
\begin{equation}
 \label{route:automorphism}
 L_1\cdots L_b=W(R_1\cdots R_b).
\end{equation}
The left side opens only positive packets. The right side opens only
negative packets and then applies one total jump. For a term at
$X_y$, the input of the total jump differs from $y$ by a multiple
of $p$. Its pairing is therefore $\Omega(p,y)$, and
Equation~\eqref{route:positive-total} makes that jump positive.

We spell out the energy in Equation~\eqref{route:automorphism}.  If the
individual increments in a term are $j_i p$, its cross-packet shift is
\begin{equation}
 \label{route:cross-energy}
 \begin{split}
 \sum_{i<j}\Omega(x_i+j_i p,x_j+j_j p)
 ={}&\sum_{i<j}\Omega(x_i,x_j)\\
 &+\sum_{i<j}\bigl(j_i\Omega(p,x_j)
                         +j_j\Omega(x_i,p)\bigr).
 \end{split}
\end{equation}
There is no quadratic term because $\Omega(p,p)=0$.  On the right
of Equation~\eqref{route:automorphism}, add the energy of the final total
jump to this expression and to the individual inverse-jump energies.
On the left, add the individual forward-jump energies.  Taking a
fixed $v$-coefficient of $X_y$ in the identity gives equality of the
two finite sets of jump tuples with exactly these energies.
Match them by rank.  Restoring the fixed underlying packet energies
adds the same integer to both sides.

The bijection is now explicit in either direction.  From the left of
Equation~\eqref{route:switch-equation}, open its positive packets by their
single-packet bijections, keep its negative packets as underlying
states, apply the jump-tuple rank bijection, and close the negative
packets into states in $I$.  Conversely, from the right, open the
negative packets, retain the positive packets, invert the jump-tuple
bijection, and close the positive packets into states in $II$.
Invariant packet types ensure that the same underlying data and rank
sets are used on reversal.  All ranks are finite by
Lemma~\ref{route:finite}.  Finally, if $\Omega(p,y)=0$, every monomial
on the line $y+\Z p$ commutes with $g_p$, so its total transform has
only the zero jump.
\end{proof}

\begin{example}[Zero total jump with changing packets]
\label{route:zero-example}
In the monomial identity of the switch, take $W=\Ad(D_p)$ and
$\Omega(p,x_1)=1$, $\Omega(p,x_2)=-1$.
The two sides of Equation~\eqref{route:automorphism} become
\[
 (X_{x_1}+X_{x_1+p})X_{x_2}
   =X_{x_1}(X_{x_2}+X_{x_2+p}).
\]
The total pairing is zero, so $W$ fixes every monomial in the
right-hand product. The terms at exponent $x_1+x_2+p$ agree because
\[
 \Omega(x_1+p,x_2)=\Omega(x_1,x_2+p)=\Omega(x_1,x_2)-1.
\]
Thus the local matching transfers the increment from the positive
packet to the negative packet while its recorded total jump is zero.
The example concerns the monomial matching with retained exponents
fixed; its point is that zero total jump need not fix each packet.
\end{example}

\subsection{Common histories and terminal anchor states}
\label{route:histories}

Before choosing a composition, fix a generic rational representative
$h_{H,y}\in H$ for every truncated chamber $H$ and every $y\in B$.
Use the line $h_{H,y}+tV(y)$, $t\geq0$.  Choose the representatives
so that its transverse wall crossings are distinct and avoid joints.
There are only finitely many avoidance conditions.  At large $t$ it
lies in a testing chamber near $V(y)$; if $V(y)=0$, the line is
constant and has no crossings.  List its crossings in the direction
of decreasing $t$, toward $H$.  Each has the sign required by
Equation~\eqref{route:positive-total}.  The choices depend on $(H,y)$ only,
not on the number, sizes, energies, or indices of packets.

For a section $P$ supported above $Nd_1$, the coefficient recursion
of Lemma~\ref{wall:recursion} is
\begin{equation}
 \label{route:history-recursion}
 [X_y]P^H=A_y(P)+
 \sum_{W:I\to II}\ \sum_{Nd_1\leq z<y}
       [X_y]W(X_z)\,[X_z]P^I.
\end{equation}
The outer sum follows the selected line, oriented toward $H$.
Only increments on the crossed ray have nonzero coefficients.
The zero-jump coefficient is one, which is why
Equation~\eqref{route:history-recursion} is a sum of increments rather than
a product of separate wall choices.

Define a total history from $(H_-,u)$ recursively.  At a pair $(H,y)$
it either stops, with terminal label $y$, or chooses one crossing
$W:I\to II$ in Equation~\eqref{route:history-recursion}, an exponent $z<y$,
and a nonzero decorated jump from $z$ to $y$, then continues from
$(I,z)$.  Its energy is the sum of its jump energies.  Let $T_x$ be
the set of histories with terminal label $x$, and $T_x(E)$ its energy
$E$ slice.  The choices of crossing and chamber at each step are
part of the history, so equal-looking jumps on different portions
of a path are distinguished.

Every nonzero jump lowers the cone height by at least one, so a
history has at most $h(R)$ such jumps.  Lemma~\ref{route:finite}
proves directly that $T_x(E)$ is finite.  Applying
Equation~\eqref{route:history-recursion} to $\Theta_x$, whose incoming
coordinates are $A_y(\Theta_x)=\mathbf1_{y=x}$, gives
\begin{equation}
 \label{route:history-series}
 \sum_{t\in T_x}v^{E(t)}=[X_u]\Theta_x^{H_-}.
\end{equation}
This is a coefficientwise identity in the Laurent completion.
The definition of $T_x$ itself uses only local monomial jumps.

The packet realization of this recursion is also reversible.  Start
with a product state of exponent $y$ in $H$, and traverse the chosen
line backwards, applying Proposition~\ref{route:switch} at each
wall.  A zero total jump continues to the next wall.  At the first
nonzero total jump, record its crossing and decoration and recurse
from the smaller exponent and its chamber.  If every total jump is
zero, stop with the product state in the testing chamber.  Between
two nonzero jumps only finitely many walls are crossed.  For the
inverse, first reconstruct the state in $(I,z)$ recursively.  Apply
the inverse local switch with the recorded jump $z\to y$ to obtain
its state in chamber II, then undo the zero-jump switches between
II and $H$.  At a stopping leaf, begin with its terminal testing
state and undo every zero-jump switch on that selected line.
Thus for every composition this procedure is a
bijection to histories with their terminal product states.  It uses
literally the same $T_x$ for all compositions.

\begin{lemma}[Return to the anchors]
\label{route:return}
A nonempty terminal product state has label
$x=x_\mu=\sum_{i=1}^D\mu_i d_i$ for a partition $\mu\vdash N$,
padded by zeroes.  Keeping its total history fixed, it can be
transported bijectively to an ordered tuple of anchor subsets of
sizes $k_1,\ldots,k_b$ and total exponent $x_\mu$.  Every such
anchor tuple has energy zero.
\end{lemma}

\begin{proof}
In ordinary commuting variables $z_1,\ldots,z_D$, expand
\begin{equation}
 \label{route:anchor-product}
 \prod_{j=1}^b e_{k_j}(z)=\sum_{\mu\vdash N}a_{\boldsymbol k,\mu}
 m_\mu(z).
\end{equation}
The nonnegative integer $a_{\boldsymbol k,\mu}$ is the number of
ordered subsets of $[D]$ of the specified sizes in which index $i$
occurs $\mu_i$ times.  By Theorem~\ref{tri:monomial}, the incoming
coordinates of $P_{\boldsymbol k}$ are exactly these constants at
$x_\mu$ and zero elsewhere.  A testing-chamber state therefore has
one of these labels, and its total coefficient there is the
constant $a_{\boldsymbol k,\mu}$.

The coefficient at the same pure label in $H_-$ is also
$a_{\boldsymbol k,\mu}$.  Indeed the initial independent-set
monomials have nonnegative vertex coordinates, so a tuple using a
bridge cannot contribute to an exponent with every bridge
coordinate zero.  All remaining tuples use anchors, which are
mutually independent and pair by zero.  Their packet energies and
cross-packet shifts are zero.

Here is a transport of the states, rather than just an equality of
their number.  Fix a generic all-negative rational covector $h_-$.
Within a sufficiently small neighborhood of the testing direction
$V(x)$, move the terminal chamber to the chamber of
$h_-+tV(x)$ for large $t$.  Only hyperplanes with
$\Omega(p,x)=0$ can separate these two perturbations.  Their
packet switches have zero total jump by
Proposition~\ref{route:switch}.  Next follow $h_-+tV(x)$ down to
$t=0$.  Every crossed positive root $p$ has
$\Omega(p,x)>0$, and the crossing is from positive to negative.
At each crossing, the coefficient of $X_x$ gains only nonnegative
contributions from lower exponents, with its old coefficient
retained by the zero jump.  Its starting and ending coefficients
are the same constant $a_{\boldsymbol k,\mu}$.  Coefficientwise
positivity forces every gain along this finite path to be zero.

It follows that, in Equation~\eqref{route:switch-equation} for these product
states at $x$, only the zero-total-jump part occurs at every wall.
The switches consequently give bijections all the way to the
literal anchor tuples in $H_-$.  Reversing the selected path gives
the inverse.  All intermediate choices can be fixed by the rational
ordering above; they may use the composition but add no total jump
to the recorded history.  The history determines its terminal
testing chamber, so the inverse loses no data.
\end{proof}

\begin{proposition}[Formal degree bounds]
\label{route:parity}
For every partition $\mu\vdash N$, histories in $T_{x_\mu}$ have
\begin{equation}
 \label{route:parity-range}
 -M\leq E(t)\leq M,\qquad E(t)\equiv M\pmod 2.
\end{equation}
Consequently $T_{x_\mu}$ is finite, and
$v^M[X_u]\Theta_{x_\mu}^{H_-}$ belongs to $\N[q]$ with $q=v^2$
and degree at most $M$.
\end{proposition}

\begin{proof}
Equation~\eqref{route:history-series} identifies the history enumerator
with $[X_u]\Theta_{x_\mu}^{H_-}$. By
Equation~\eqref{tri:positive-coefficient}, this coefficient is
$v^{-M}c_\mu(v^2)$, where $c_\mu(q)\in\N[q]$ has degree at most $M$.
Its allowed powers of $v$ are therefore exactly among
$-M,-M+2,\ldots,M$. Histories have nonnegative multiplicities, so
none can occur at any other energy. Each of these finitely many
energy slices is finite by Lemma~\ref{route:finite}; hence their
union $T_{x_\mu}$ is finite as well. Multiplication by $v^M$ changes
each history weight to $q^{(E(t)+M)/2}$, proving the assertion.
\end{proof}

\subsection{Required masks and exact masks}
\label{route:masks}

For a permutation word $w=w_1\cdots w_N$ on the indicated vertices,
define its nonedge ascent mask by
\[
 \mathcal P_G(w)=
 \{j: w_j<w_{j+1},\ \{w_j,w_{j+1}\}\notin E(G)\}.
\]
For a word $s=s_1\cdots s_N$ on the labeled anchor indices
$[D]$, define its ascent mask by
$\mathcal Q(s)=\{j:s_j<s_{j+1}\}$.
For a subset $J\subseteq[N-1]$, a \emph{required} mask means
$J\subseteq\mathcal P_G(w)$ or $J\subseteq\mathcal Q(s)$;
an \emph{exact} mask means equality.  These are different conditions.

Place packet boundaries exactly at the complement of $J$.  This
gives a composition ${\boldsymbol k}(J)$ of $N$.  If
$J\subseteq\mathcal P_G(w)$, each packet is an increasing independent
set.  In fact, if successive letters $a<b$ in a packet satisfy
$h(a)<b$, then every later letter $c\geq b$ also satisfies
$h(a)<c$.  Hence every pair in that packet is a nonedge.
Conversely, an independent set in increasing order has precisely
the required internal ascents.  Product states at $u$ for this
composition are therefore exactly the words with the required
mask $J$.  Their energies are
\begin{equation}
 \label{route:word-energy}
 E(w)=2\inv_G(w)-M,
\end{equation}
since no edge is internal to a packet and
Equation~\eqref{route:product-energy} counts every edge between packets.

At a terminal label $x_\mu$, writing each anchor subset in increasing
index order gives a word with multiplicities
$(\mu_1,\ldots,\mu_D)$ and with $J\subseteq\mathcal Q(s)$.
Lemma~\ref{route:return} and the history construction therefore give
a computable bijection, with computable inverse,
\begin{equation}
 \label{route:required-bijection}
 B_J:\{w:J\subseteq\mathcal P_G(w)\}
 \longrightarrow
 \coprod_{\mu\vdash N}
 \{(t,s):t\in T_{x_\mu},\
       \operatorname{mult}(s)=\mu,\ J\subseteq\mathcal Q(s)\}.
\end{equation}
It satisfies $E(t)=E(w)$.  The target histories in this formula are
the same for every $J$; only the permitted words $s$ change.

We convert the required-mask bijections to exact-mask bijections by
finite bijective cancellation, in the spirit of the Garsia--Milne
involution principle \cite{GM81}.  The complete alternating-path
argument below also supplies termination in both directions.

\begin{lemma}[Exact-mask refinement]
\label{route:exact}
For every $J\subseteq[N-1]$ there is a computable energy-preserving
bijection $\Phi_J$ between the two sides of
Equation~\eqref{route:required-bijection} with their masks required to be
\emph{exactly} $J$.  Both directions terminate on every input.
\end{lemma}

\begin{proof}
Construct the bijections in decreasing order of $|J|$, breaking ties
by the fixed subset order.  For an exact-$J$ input $w$, apply $B_J$.
If its output $(t,s)$ has exact mask $J$, stop.  Otherwise put
$K=\mathcal Q(s)\supsetneq J$, apply the already constructed inverse
$\Phi_K^{-1}$, and apply $B_J$ again.  Repeat this instruction.
For the inverse procedure start at an exact-$J$ pair, apply
$B_J^{-1}$, and, whenever the resulting word has exact mask
$K\supsetneq J$, apply $\Phi_K$ followed by $B_J^{-1}$.
Thus every nested call involves a strictly larger mask.

To prove that the outer repetition terminates, fix its energy $E$.
The required-$J$ source is finite, and its target is finite by
Lemma~\ref{route:finite} (also by
Proposition~\ref{route:parity}).  On their disjoint union draw the
matching given by $B_J$.  On the vertices whose exact masks strictly
contain $J$, draw a second matching, the disjoint union of the
already defined $\Phi_K$ for $K\supsetneq J$.  Each vertex has one
edge of the first matching, and either zero or one edge of the
second.  Distinct edges that happen to have the same endpoints are
retained as distinct matching edges.
Figure~\ref{route:matching-figure} illustrates the alternating path in
their union.

An exact-$J$ starting vertex is an endpoint of this finite graph.
Its component is a path, not a cycle: all vertices have degree at
most two, and a component containing a degree-one vertex cannot be
a cycle.  Following its alternating edges therefore never revisits
a vertex and eventually reaches the other endpoint.  A return edge
of the second matching always lands at a source with a larger
exact mask, where the first matching can be followed again.
Consequently the other endpoint lies on the target side and has
exact mask $J$.  Reversing the same path is precisely the inverse
procedure described above.  Every edge preserves $E$, proving the
claimed energy preservation as well as termination.
\end{proof}

\begin{figure}[htbp]
\centering
\begin{tikzpicture}[>=Latex,
  vertex/.style={draw,rounded corners,inner sep=4pt,font=\small}]
 \node[vertex] (l0) at (0,1.3) {$w:\ J$};
 \node[vertex] (r1) at (2,0) {$(t_1,s_1):K_1$};
 \node[vertex] (l1) at (4,1.3) {$w_1:K_1$};
 \node[vertex] (r2) at (6,0) {$(t_2,s_2):K_2$};
 \node[vertex] (l2) at (8,1.3) {$w_2:K_2$};
 \node[vertex] (r0) at (10,0) {$(t,s):J$};
 \draw[->,thick] (l0)--node[above,sloped,font=\small] {$B_J$} (r1);
 \draw[->,thick,dashed] (r1)--node[above,sloped,font=\small]
   {$\Phi_{K_1}^{-1}$} (l1);
 \draw[->,thick] (l1)--node[above,sloped,font=\small] {$B_J$} (r2);
 \draw[->,thick,dashed] (r2)--node[above,sloped,font=\small]
   {$\Phi_{K_2}^{-1}$} (l2);
 \draw[->,thick] (l2)--node[above,sloped,font=\small] {$B_J$} (r0);
 \node[font=\small,align=center] at (5,-0.9)
   {$K_i\supsetneq J$ at every interior pair; endpoints have exact mask $J$.\\
    Solid: required-mask matching.\quad Dashed: larger exact-mask matching.};
\end{tikzpicture}
\caption{The exact-mask algorithm follows a finite alternating path.
An endpoint cannot enter a cycle in the union of two matchings.}
\label{route:matching-figure}
\end{figure}

\subsection{The partition assigned to a word}
\label{route:algorithm}

Let $\mathcal A_G=\{w:\mathcal P_G(w)=\varnothing\}$ be the
admissible domain.  For $w\in\mathcal A_G$, compute
\[
 \Phi_{\varnothing}(w)=(t,s),\qquad t\in T_{x_\mu},
 \qquad \Lambda_G(w)=\mu.
\]
The finite algorithm on an individual input has four steps.
\begin{enumerate}
\item Construct the triangular embedding and finite cone interval, and
fix the common rational path representatives for every required $(H,y)$
before choosing any packet composition.
\item For each required mask encountered, form its packets and evaluate
$B_J$ through the local packet switches, common histories, and return
to anchor subsets. Positive and negative packets use the two directions
of the local switch, although the total jump has the positive sign.
\item Evaluate $\Phi_{\varnothing}$ by the exact-mask recursion in
Lemma~\ref{route:exact}, following its alternating walk until the output
has exact empty mask.
\item Read the multiplicities of the resulting weakly decreasing anchor
word as the partition $\Lambda_G(w)$.
\end{enumerate}
Nested calls strictly increase mask size, while each outer repetition
follows an endpoint path in the finite union of two matchings. Thus the
call depth is bounded and each alternating walk terminates. There is
no need to enumerate the final fibers of $\Lambda_G$ or to supply the
final elementary coefficients. Section~\ref{main:completion} gives the
reversal for the original input convention.
If the input is supplied in the insertion digits of
Lemma~\ref{graph:digits}, first decode those digits by largest-vertex
insertions.  The construction then uses the resulting admissible
word.  Empty input is assigned the empty partition.

\begin{theorem}[Permutation witness]
\label{route:witness}
For every naturally labeled unit interval graph $G$, the preceding
finite rule assigns a partition $\Lambda_G(w)\vdash n$ to each
$w\in\mathcal A_G$ and satisfies the formal polynomial identity
\begin{equation}
 \label{route:witness-identity}
 X_G(X;q)=\sum_{w\in\mathcal A_G}
                   q^{\inv_G(w)}e_{\Lambda_G(w)}(X).
\end{equation}
In particular, every elementary coefficient is in $\N[q]$ and has
degree at most $|E(G)|$.
\end{theorem}

\begin{proof}
For any padded partition $\mu$, exactly one word of multiplicities
$(\mu_1,\ldots,\mu_D)$ has no strict ascent: the weakly decreasing
word.  This remains true when parts of $\mu$ repeat, since the
anchor indices themselves are distinct labels.  Consequently
$\Phi_{\varnothing}$ identifies the individual words of shape
$\mu$ with $T_{x_\mu}$.  Energy preservation,
Equation~\eqref{route:word-energy}, and Equation~\eqref{route:history-series} give
\[
 \sum_{\substack{w\in\mathcal A_G\\\Lambda_G(w)=\mu}}
       v^{2\inv_G(w)-M}
   =\sum_{t\in T_{x_\mu}}v^{E(t)}
   =[X_u]\Theta_{x_\mu}^{H_-}
   =[X_u]m_\mu(Y).
\]
Multiply by $v^M$.  The conversion $q=v^2$ is legitimate as a
formal polynomial conversion by
Proposition~\ref{route:parity}, including both its parity and its
degree bounds.  The coefficient identity of
Proposition~\ref{graph:kernel} now identifies the right side with
$[e_\mu]X_G(X;q)$.  This proves the identity for each shape and
hence Equation~\eqref{route:witness-identity}. The algorithms producing
$\Lambda_G(w)$ use only elementary packets and local wall jumps;
the coefficient identity certifies their output.
\end{proof}
